\section{Output events and exclusion bases}
\label{app:bases}

This appendix first gives the formal definition of the lifted output events of
Section~\ref{sec:lifted-output-events}.  It then supplies the details behind
Section~\ref{sec:systems}: the canonical exclusions of
Sections~\ref{sec:functional-calm}--\ref{sec:complete-calm}
are exclusions (Definition~\ref{def:exclusion}) and form a canonical basis,
compound exclusions behave like single ones, and every run under any basis
has an observably equivalent canonical run that retains at least as many rows
after each event.  Exclusions are partial functions, so a basis may make an
exclusion available only at states satisfying a precondition of its own.

\subsection{Lifted output events}
\label{app:lifted-output-events}

Fix a state $(i,R)$ and an output $o\in R(i)$.  A row is \emph{ruled out by}
$o$ when its output does not \outrelbase $o$.  The ruled-out rows fall into two
disjoint sets:
\begin{align}
  \Subsumed_o&=\{(k,x)\in R\setminus R[\mathrm{out}\OutputGe o]:\notag\\
    &\qquad\quad x\OutputOrd u\text{ for some }u\in R[\mathrm{out}\OutputGe o](k)\},
      \label{eq:output-subsumptions}\\
  \Excluded_o&=(R\setminus R[\mathrm{out}\OutputGe o])\setminus\Subsumed_o.
      \label{eq:exclusion-removals}
\end{align}

\begin{definition}[Lifted output event, output subsumption]
\label{def:lifted-output-event}
Let $(i,R)$ be a state and $o\in R(i)$.  The \emph{output event}
$\mathsf{output}(o)$ keeps exactly the rows whose outputs \outrelbase $o$ and
emits $o$, atomically, ending at $(i,R[\mathrm{out}\OutputGe o])$ as in
Equation~\ref{eq:output-action}.  It removes the rows in $\Subsumed_o$, each
\outrelpast by a kept row at the same input, by \emph{output subsumption},
which is not an exclusion.
\end{definition}

\begin{definition}[Canonical exclusion]
\label{def:canonical-exclusion}
When $\Excluded_o\neq\varnothing$, the output event applies the \emph{canonical
exclusion} $x_o$:
\begin{equation}
  x_o(i,R)=\bigl(i,\ R\setminus\Excluded_o\bigr).
  \label{eq:canonical-exclusion-relational}
\end{equation}
At an input $k$, it excludes input $k$ if it removes every row there; otherwise
it excludes each output whose row it removes at $k$.
\end{definition}

In Figure~\ref{fig:ambiguity-frames}(f), emitting $[A,B]$ places the row for
$[B,A]$ in $\Excluded_o$, so the event excludes that output at $\{A,B\}$.
In panel~(g), the two removed route sets are in $\Subsumed_o$, each
\outrelpast by a retained route set, so the event subsumes them without
exclusion.  With one row per input, $\Subsumed_o$ is empty, and $x_o$ excludes
each input whose row is removed; this is the functional output event of
Definition~\ref{def:output-event}.  Apart from these, a row leaves the
remainder table only by input loss (\S\ref{sec:functional-events}).

Legality lifts with the event: $o$ is a legal output when $\mathsf{output}(o)$
applies no exclusion, exactly when $\Excluded_o$ is empty, though the event may
still subsume the rows in $\Subsumed_o$.  The legal-output test lifts as well:

\begin{proposition}[Legal-output test]
\label{prop:exposability}
An output $o\in R(i)$ is legal at $(i,R)$ if and only if, for every legal
$j\InputGe i$ and $x\in R(j)$, some $u\in R(j)$ \outrel both:
\begin{equation}
  o\OutputOrd u
  \qquad\text{and}\qquad
  x\OutputOrd u.
  \label{eq:safe-exposure}
\end{equation}
\end{proposition}

\begin{proof}
The event $\mathsf{output}(o)$ applies no exclusion exactly when every
ruled-out row $(j,x)$ lies in $\Subsumed_o$, that is, is \outrelpast by some row
$(j,u)$ \outrelpart $o$.  Rows that already \outrelbase $o$ satisfy the stated
test with $u=x$, so this is exactly the stated test.
\end{proof}

This test marks the exclusion boundary: if it fails, emitting $o$ requires an
exclusion.

\subsection{Exclusion bases and compound exclusions}
\label{app:exclusion-bases}

\begin{definition}[Canonical basis]
\label{def:canonical-exclusion-basis}
The \emph{canonical exclusion basis} collects the canonical exclusions of
Definition~\ref{def:canonical-exclusion},
\begin{equation}
  \mathcal B_P
    =\{x_o:o\in O,\ \operatorname{dom}(x_o)\neq\varnothing\},
  \label{eq:canonical-exclusion-basis-relational}
\end{equation}
where $x_o$ is defined at the states with $o\in R(i)$ and $\Excluded_o\neq\varnothing$.
\end{definition}

\begin{definition}[Canonical event]
\label{def:canonical-event}
At $(i,R)$, let $\kappa_o=\varepsilon$ when $\Excluded_o=\varnothing$, and
$\kappa_o=[x_o]$ otherwise.  The \emph{canonical event} for $o$ is
$\Out{\kappa_o}{o}$: the output event of
Definition~\ref{def:lifted-output-event}, and with one row per input that of
Definition~\ref{def:output-event}.
\end{definition}

\paragraph{The canonical exclusion is an exclusion.}
Suppose a removed row $(k,x)\in\Excluded_o$ were \outrelpast by a row $(k,y)$
that $x_o$ leaves.  Either $y$ \outrel $o$, or $y\in\Subsumed_o$ and is
\outrelpast by a row \outrelpart $o$.  By transitivity, $x$ is then
\outrelpast by a row \outrelpart $o$, so $x\in\Subsumed_o$, a contradiction.

\begin{lemma}[Containment, canonical support, and exclusion-free events]
\label{lem:canonical-exclusion-normal-form}
Let $(i,R)$ be a state and $o\in R(i)$.
\begin{enumerate}
  \item Every legal output event $\Out{\gamma}{o}$ takes $(i,R)$ to some $(i,R')$ with
        $R'\subseteq R[\mathrm{out}\OutputGe o]$, and removes the rows in
        $\Excluded_o$ by exclusion.
  \item The canonical event $\Out{\kappa_o}{o}$ is legal and takes $(i,R)$ to
        $(i,R[\mathrm{out}\OutputGe o])$, so it retains every row that any legal
        $\Out{\gamma}{o}$ retains.
  \item If $\Out{\varepsilon}{o}$ is legal, it takes $(i,R)$
        to $(i,R[\mathrm{out}\OutputGe o])$; it is legal exactly when $o$ is
        a legal output.
\end{enumerate}
\end{lemma}

\begin{proof}
(1) A legal event ends at $R_c[\mathrm{out}\OutputGe o]\subseteq
R[\mathrm{out}\OutputGe o]$, since its exclusions only remove rows.  A row in
$\Excluded_o$ has no \outrelpart row in $R[\mathrm{out}\OutputGe
o](k)$, and hence none in $R_c[\mathrm{out}\OutputGe o](k)$; it cannot be
subsumed, so a legal event must remove it by exclusion.  (2) The canonical
exclusion removes exactly the rows in $\Excluded_o$, all
of them ruled out by $o$, so $R_c[\mathrm{out}\OutputGe o]=R[\mathrm{out}\OutputGe
o]$, which contains $(i,o)$.  Every remaining ruled-out row is in $\Subsumed_o$, and its
\outrelpart row is not in $\Excluded_o$, so it
survives the exclusion.  The event is therefore legal and ends at
$R[\mathrm{out}\OutputGe o]$; by (1) no legal event keeps more.  (3) With no
exclusion, $R_c=R$, and a legal event ends at $R[\mathrm{out}\OutputGe o]$.
It is legal exactly when every ruled-out row is in $\Subsumed_o$, that is, when $\Excluded_o$ is empty, which makes $o$ a legal output.
\end{proof}

\begin{definition}[Canonical system]
\label{def:canonical-system}
A system (Definition~\ref{def:system}) is \emph{canonical} when its basis is
$\mathcal B_P$ and each output event it performs is the canonical event $\Out{\kappa_o}{o}$ at the
participant's state where it occurs, and it performs no exclusion event; its histories
are then exactly those of Section~\ref{sec:histories}.
\end{definition}

\begin{lemma}[Compound exclusions are exclusions]
\label{lem:compound-exclusion}
The composite of a nonempty sequence of exclusions, defined wherever the
sequence is, is an exclusion.
\end{lemma}

\begin{proof}
Each step leaves the current input unchanged and maps a state to a state, so the
composite does too, and it removes a nonempty set of rows because the first
step does.  Suppose step $m$ removes $(k,u)$, and let $R_m$ be the table that
step $m$ leaves.  By Equation~\ref{eq:exclusion-excludes} for that step, no
row of $R_m$ at $k$ \outrel $u$; the final table is contained in $R_m$, so no
remaining row at $k$ \outrel $u$ either.
\end{proof}

A compound output event (Definition~\ref{def:system-output-event}), or a
compound exclusion event, therefore has the same effect as an event that applies the composite exclusion alone,
and Lemma~\ref{lem:canonical-exclusion-normal-form} covers it as stated.  In a
functional specification the composite removes every row at each input it
touches, so it empties.

\begin{definition}[Run under a basis]
\label{def:run-under-basis}
Fix an exclusion basis $\mathcal B$.  A \emph{run under $\mathcal B$} is a
finite transition sequence from a canonical state $\sigma_P(i)$ of input events, legal output events $\Out{\gamma}{o}$
(Definition~\ref{def:system-output-event}), and exclusion events
(Definition~\ref{def:exclusion-event}) that use only exclusions from
$\mathcal B$.
\end{definition}

\begin{definition}[Canonical run]
\label{def:canonical-run}
The \emph{canonical runs} are the runs under the canonical basis
$\mathcal B_P$ of Definition~\ref{def:canonical-exclusion-basis} whose output
events are all canonical; they are the runs of
Sections~\ref{sec:functional-calm}--\ref{sec:complete-calm}.
\end{definition}

To make the claim of Lemma~\ref{lem:canonical-equivalent} precise, call two
runs \emph{observably equivalent} when they start at the same canonical state and agree once exclusion sequences are erased and exclusion events dropped:
wherever one has $\mathsf{input}(j)$ the other does, and wherever one has
$\Out{\gamma}{o}$ the other has $\Out{\gamma'}{o}$ for some $\gamma'$.  They accept the same inputs and emit
the same outputs, in the same order, which is all the environment observes.
Observably equivalent runs may differ in the exclusions their output and exclusion events apply, and
hence in their remainder tables.

\begin{proof}[Proof of Lemma~\ref{lem:canonical-equivalent}]
By induction on the events.  Both runs start at the same canonical state, so
their remainder tables are initially equal.  Suppose that before some event the
original run's table $R$ and the canonical run's table $\widehat R$ satisfy
$R\subseteq\widehat R$; we show that the inclusion holds afterward.  An input
event $\mathsf{input}(j)$ that is legal at $R$ is also legal at $\widehat R$.
It causes the table to lose every row whose input does not \inrelbase $j$, keeping
$R[\mathrm{in}\InputGe j]$ in the original run and
$\widehat R[\mathrm{in}\InputGe j]$ in the canonical run.  Since
$R\subseteq\widehat R$, the first remains contained in the second.  For a legal
output event $\Out{\gamma}{o}$ with $o\in R(i)$ taking $R$ to $R'$, containment
(Lemma~\ref{lem:canonical-exclusion-normal-form}(1)) gives
$R'\subseteq R[\mathrm{out}\OutputGe o]$.  Since
$o\in R(i)\subseteq\widehat R(i)$, canonical support
(Lemma~\ref{lem:canonical-exclusion-normal-form}(2)) applies at $\widehat R$:
$\Out{\kappa_o}{o}$, with $\kappa_o$ taken at $(i,\widehat R)$, is legal and
takes $\widehat R$ to $\widehat R[\mathrm{out}\OutputGe o]$, which contains
$R[\mathrm{out}\OutputGe o]$.  An exclusion event keeps the current input and
only removes rows, so the canonical run skips it and the inclusion still
holds.
\end{proof}

\begin{corollary}[Reconciliation under any basis]
\label{cor:reconciliation-any-basis}
Let runs under any exclusion bases end at $(i_r,R_r)$, and let $j$ be a
common \indesc of the $i_r$.  If some continuations, whose output and exclusion events
may apply any exclusions, bring them to a common state with current input $j$,
then some $u\in\fiber{j}$ \outrel every output emitted on every run.
\end{corollary}

\begin{proof}
The common state has a row $(j,u)$, and events never add rows, so
$(j,u)\in R_r$.  By Lemma~\ref{lem:canonical-equivalent}, $R_r$ is contained in
the table of a canonical run with the same inputs and outputs, and
Lemma~\ref{lem:run-residual} makes $u$ \outrel the outputs of run $r$.
\end{proof}

The converse of Lemma~\ref{lem:canonical-equivalent} fails: a canonical run
need not be observably equivalent to any run under a particular system's
basis.  That basis may lack the exclusion an
output event requires, and so disable the output altogether.  Restricting
attention to one system's runs could therefore hide non-monotonicity.  This is
why consistency (Definition~\ref{def:consistency}) is defined over canonical
runs rather than over a system's own.

\paragraph{Application-specific exclusions.}
Locks and timestamps are exclusions of this kind: two-phase locking applies an exclusion that fixes an order among some transactions, and timestamp assignment places one transaction
relative to all others.  Section~\ref{sec:systems} illustrates both, and
Appendix~\ref{app:cc-bases} writes the two bases down.
